\section{The long-form proof format}\label{proofs: sec: long form}\doHeader{Long-form proof format}

The long-form proof format uses a particular syntactical form
that helps make the details and logical structure of the proof particularly clear.
This section defines the format and gives examples.
Section~\ref{proofs: sec: exercises with solutions} has more examples.

\subsection{There are four kinds of proof steps}
\begin{figure}
  \begin{framed}\small
    \begin{theorem}\label{proofs: thm: first}
      The number  $x =\log_9 12$ is irrational. 
      (That is, the $x$ such that $9^x = 12$.)
    \end{theorem}
    \begin{longFormProof}
      \step Assume for contradiction that $x$ is rational.
      \comment{\ldots starts a proof-by-contradiction block}
      \step Let $p$ and $q>0$ be integers such that $p/q {} = \log_9 12$.
      \comment{\ldots definition}
      \step Raising 9 to the power of each side, it must be that $9^{p/q} = 12$.
      \comment{\ldots statement}
      \step\label{proofs: step: eqn}
      Raising each side to the power of $q$, it must be that $9^p = 12^q$.
      \comment{\ldots statement}
      \step Because $q>0$, the right-hand side $12^q$ is a multiple of 2.
      \comment{\ldots statement}
      \step But the left-hand side $9^p$ cannot be a multiple of 2.  So $9^p \ne 12^q$.
      \comment{\ldots statement}
      \step\label{proofs: step: contr}
      This contradicts Line~\ref{proofs: step: eqn},  so our assumption in Line~1 is wrong.  So $x$ is irrational.
      \comment{\ldots statement}
    \end{longFormProof}
    \vspace*{-10pt}
  \end{framed}
  \vspace*{-10pt}
  \caption{A long-form proof with one \emph{proof by contradiction} block.}\label{proofs: fig: contradiction}
  \vspace*{-5pt}
\end{figure}
Each long-form proof is given as a sequence of numbered steps.
Each step takes up about one line on the page.  Figure~\ref{proofs: fig: contradiction} shows an example. 
There are just four kinds of steps:
\begin{enumerate}[itemsep=0pt]
\item A \emph{definition} of a term or variable (for example, Step 2 in the figure defines $p$ and $q$),
\item A \emph{statement of fact}, with a justification (Steps 3--7), 
\item A \emph{comment} to guide the reader (the proof above has none),
\item The first line of a \emph{block} (such as in Step 1); per Section~\ref{proofs: sec: blocks},
  blocks are used to manage assumptions.   Each block takes one of four forms:
  \begin{enumerate}[label=(\roman*)]
  \item a proof by contradiction,
  \item a proof of an ``if then'' statement,
  \item a proof by cases,
  \item a proof of a ``for all'' statement.
  \end{enumerate}
\item Some combination of the above (e.g.~a definition and a statement; do this rarely, though).
\end{enumerate}

A \textbf{definition} gives a name to some quantity (beware: it must be clear that the quantity does exist).
The definition should be mathematically precise, with no ambiguity.
The definition holds throughout the remainder of the block in which it occurs.
(Just as declaring a variable in a computer program makes the variable available
for the rest of the code block in which the declaration occurs.)

Each \textbf{statement of fact} states some fact, along with a short justification.
\emph{The given fact must follow directly from the facts established in previous steps},
and your anticipated reader should be able to easily see why this is the case.
The level of detail you give thus depends on your anticipated reader.
For example, the proof in Figure~\ref{proofs: fig: contradiction} assumes the reader is familiar with the definitions of ``rational''
and ``irrational'', and with the basic properties of integers,
and it uses such properties without explaining why they hold.
When writing up exercises for these notes, you should assume the reader is another student
who knows the material in the notes,
but who doesn't yet know anything about the line of reasoning you have in mind.

When you are starting out, if you're not sure how much detail to give, err on the side
of giving more.  For the reader, it's easier to skip unnecessary detail than to figure out what
you might have had in mind when details are missing.  And gaps (missing details)
often hide errors.

As necessary, the statement of fact is accompanied by a justification to guide the reader,
along with explicit references to any previous steps that it depends on.
(By default, each step depends on the previous step,
so if that's the only dependence you don't have to mention it.)

\subsection{There are four kinds of blocks}\label{proofs: sec: blocks}
Long-form proofs can have multiple, possibly nested \textbf{blocks}.
(The proofs in Figures~\ref{proofs: fig: contradiction} and~\ref{proofs: fig: if then} each have just one, though,
the main outer block.)
Blocks are used to keep track of assumptions made during the proof.
\emph{The first line of each block introduces an assumption} (in some cases only implicitly).
The assumption is then in effect until the end of the block.
%
That is, within the block we use the assumption as if it were a true fact.
We don't need to know whether the assumption is actually true.
Rather, within the block we are imagining that it is true,
and exploring what the logical consequences would be if it were.
This kind of exploration lets us figure out facts that hold even without any assumptions.

Blocks can be used in long-form proofs in (only) the following four specific ways:

\paragraph{1. Blocks for proof by contradiction.}
The proof in Figure~\ref{proofs: fig: contradiction} has just one block---its main block, consisting of Steps~1--\ref{proofs: step: contr}.
The assumption that starts the block is that the given quantity $\log_9 12$ is rational.
The block explores the logical consequences of that assumption,
showing that, if the assumption is true,
then \emph{$9^p = 12^q$ and $9^p \ne 12^q$ (for some $p$ and $q$).}
That is, a statement and its negation both hold.
This is called a \emph{contradiction}.
Since it cannot be true, we deduce that our original assumption (which implies it) must be false.
That is, $\log_9 12$ cannot be rational.  It must be irrational.

This is called \emph{proof by contradiction}.  To do proof
by contradiction, use a block whose first line explicitly assumes the opposite
of what you want to prove.  To signal the reader,
in this first line, explicitly use a phrase including the word ``contradiction'' such as 
\emph{``Assume for contradiction that\ldots''}
Then, in the body of the block, somehow deduce a contradiction---a statement and its negation.
If you can do that, you can conclude that the assumption that started the block must be false,
and the negation of that assumption must be true.
(By the way, proof by contradiction is usually just a convenient way to organize a proof.
That is, it is usually possible to give an equivalent ``direct'' proof without using proof by contradiction.)

% Often, the first step in proving a statement
% is to identify the logical structure of the statement.
% The logical structure of the statement will guide the structure of the proof.
% The statements we'll want to prove in this class
% are, generally, expressible in first-order logic,
% using just existential and universal quantifiers ($\exists$, $\forall$)
% and logical operators (and, or, not, if/implies, if and only if).

% This brings us to the subject of \emph{blocks} within the proof.
% The previous example used just a single outer block.
% But blocks can also nest.
% That is, blocks can have sub-blocks.


\begin{figure}
\begin{framed} \small
  \begin{theorem}
    Let $n = 31415923576480135764801$.  If $\log_7 n$ is rational, then it is an integer.
  \end{theorem}
  \begin{longFormProof}
    \step Assume that $\log_7 n$ is rational.
    \comment{\ldots starts an ``if then'' block}
    \step Let $p$ and $q$ be integers with $q>0$ such that $p/q = \log_7 n$.
    \comment{\ldots definition}
    \step Raising 7 to the power of each side gives $7^{p/q} {} = n$.  \comment{\ldots statement}
    \step Raising each side to the power of $q$ gives $7^p {} = n^q$.\label{proofs: PI: EQ}
    \step As each integer has a unique prime factorization, this implies that $n$ is a power of 7.
    \comment{\smash{$\vdots$}~~~}
    \step That is, $n=7^i$ for some integer $i > 0$.
    \step Substituting that into the equation in Step~\ref{proofs: PI: EQ}, by algebra we have
    \begin{align*}
      7^p & = n^q = (7^i)^q = 7^{iq} \\
      p & = iq & \comment{(taking the $\log_7$ of each side)} \\
      p/q & = i & \comment{(dividing each side by $q>0$)} \\
      \log_7 n & = i & \comment{(using $\log_7 n = p/q$).}
    \end{align*}
    \step It follows that $\log_7 n$ is an integer (if it is rational).
  \end{longFormProof}
  \vspace*{-10pt}
\end{framed}
\vspace*{-10pt}
\caption{A long-form proof with one \emph{if-then} block.}\label{proofs: fig: if then}
\vspace*{-5pt}
\end{figure}

\paragraph{2. Blocks to prove ``if then'' statements.}
The logical meaning of a statement of the form ``if A then B'' 
is that, if A is true, then B must also be true.
You can prove such a statement without knowing whether or not A is true.
Start a block by assuming A,
then, under that assumption, show that B must also hold.
If you can do that, then B is a logical consequence of A,
so, if A is true, then B must also be true.
That is, ``if A then B''.

Figure~\ref{proofs: fig: if then} shows an example.
For some particular $n$, it shows that
if $\log_7 n$ is rational, then $\log_7 n$ is an integer.
It doesn't prove whether $\log_7 n$ is rational or not.
Rather, it shows that, if we assume that $\log_7 n$ is rational,
it follows by logic and algebra that $\log_7 n$ must in fact be an integer.

\paragraph{3. Blocks to prove ``for all'' statements.}
A ``for all'' statement is a statement of the following form,
given some set $N$ and predicate $Q$:
\begin{quote}
``for all $n$ in $N$, the predicate $Q(n)$ holds''.
\end{quote}

To prove such a statement, use a ``for all'' block.
Start the block with a line such as~\footnote
{For ``for all'' blocks, the underlying assumption (not explicitly stated)
  is that $N$ is non-empty, so $n$ is well-defined.}
\begin{quote}
  ``\emph{\ldots 1. Consider an arbitrary $n$ in $N$.}'' 
\end{quote}
This asks the reader to consider a generic element of the set $N$,
referred to by the name $n$.
To let the reader know that your are starting a ``for all'' block,
make sure you include the qualifier ``arbitrary'' in the first step.
Within the block,
assuming only that $n$ is in $N$ (and knowing nothing else about $n$)
prove that $Q(n)$ must hold.
Suppose you can indeed prove that.
Because your reasoning could be applied generically to any particular element of $N$,
it must be that $Q(n)$ holds for every particular $n$.
That is, you can deduce the following statement (outside the block!):
\emph{for all $n$ in $N$, the predicate $Q(n)$ holds.}

For example, note that the proof in Figure~\ref{proofs: fig: if then} doesn't use the specific properties
of the particular integer $n$ defined there.
So, we can apply the reasoning in that proof to \emph{any} particular positive integer $n$,
to show that if $\log_7 n$ is rational, then it is an integer.
In other words,
\emph{for all positive integers $n$, if $\log_7 n$ is rational, then it is an integer.}
\begin{figure}
\begin{framed}
  \begin{theorem}
    For all positive integers $n$, if $\log_7 n$ is rational, then it is an integer.
  \end{theorem}
  \begin{longFormProof}
    \step[proofs: A0] Consider an arbitrary positive integer $n$.
    \comment{\ldots starts a ``for all'' block}
    \begin{block}\label{proofs: MCS17: B}
      \step Assume that $\log_7 n$ is rational.
      \comment{\ldots starts an ``if then'' block}
      \step Let $p$ and $q$ be integers with $q>0$ such that $p/q = \log_7 n$.
      \comment{\ldots definition}
      \step Raising 7 to the power of each side gives $7^{p/q} {} = n$.  \comment{\ldots statement}
      \step Raising each side to the power of $q$ gives $7^p {} = n^q$.\label{proofs: PI2: EQ}
      \step As each integer has a unique prime factorization, this implies that $n$ is a power of 7.
      \comment{\smash{$\vdots$}~~~}
      \step That is, $n=7^i$ for some integer $i > 0$.
      \step Substituting that into the equation in Step~\ref{proofs: PI2: EQ}, by algebra we have
      \begin{align*}
        7^p & = n^q = (7^i)^q = 7^{iq} \\
        p & = iq & \comment{(taking the $\log_7$ of each side)} \\
        p/q & = i & \comment{(dividing each side by $q>0$)} \\
        \log_7 n & = i & \comment{(using $\log_7 n = p/q$).}
      \end{align*}
      \step It follows that $\log_7 n$ is an integer.
    \end{block}
    \step[proofs: AN] By Block~\ref{proofs: MCS17: B}, if $\log_7 n$ is rational, then $\log_7 n$ is an integer.
  \end{longFormProof}
  \vspace*{-10pt}
\end{framed}
\vspace*{-10pt}
\caption{This proof has a \emph{for-all} block containing an \emph{if-then} block.}\label{proofs: fig: for all}
\end{figure}
Figure~\ref{proofs: fig: for all} shows the corresponding long-form proof.
It uses a ``for all'' block as described above.
The logical structure of the proof
mirrors the logical structure of the statement:
the main, outer block (Steps~\ref{proofs: A0} through~\ref{proofs: AN}) is a ``for all'' block;
the inner block (Block~2) is an ``if then'' block.

\paragraph{4. Blocks for proof by cases.}
The fourth and final type of block is used for case analysis.
Suppose you know that
``\emph{$C_1 \text{ or } C_2$}'' holds,
and you want to show that some proposition $X$ holds.
(You may not know which of $C_1$ or $C_2$ holds,
but you know at least one has to.)

First, consider the case that $C_1$ holds.
Start a block with the assumption that $C_1$ holds.
Under that assumption (within that block) show somehow that $X$ holds.
Then you know that, if $C_1$ holds, then $X$ holds.

Then consider the other case, that $C_2$ holds.
Start a block with the assumption that $C_2$ holds.
Under that assumption (within that block), show somehow that $X$ holds.
Then you know that, if $C_2$ holds, then $X$ holds.

Finally, since you know that $C_1 \text{ or } C_2$ holds,
and you know that in either case $X$ must hold,
you can conclude (without any assumptions) that $X$ holds.

This is the principle of proof by cases.
It extends naturally to situations with more than two cases,
as long as you know that at least one of case's assumption must hold.
%
When you start a block for a proof by cases,
signal it by starting the statement with a label, such as ``\underline{Case 1.} \ldots''.
You can also use a phrase such as ``\emph{Consider the case} that \ldots'' to signal the reader.
If necessary, use a comment to the reader before the cases
to help the reader understand their structure.

Figure~\ref{proofs: fig: cases} shows an example.
Here's the context for the theorem shown there.
Assume that, for every pair of people in the world, either the two have met, or they have not.
For any given group $G$ of people,
if every pair in $G$ has met, call $G$ a \emph{club}.
If every pair in $G$ has not met, call $G$ a group of \emph{strangers}.
The theorem says that in every group of six,
there must be either a group of three who form a club,
or a group of three strangers.
The proof uses two cases, each with two subcases.

It is interesting to note that in that proof,
the part that proves that $G$ is good is constructive.
That is, it yields an efficient algorithm, which, given any group $G$ of six people,
finds three that are a club or are strangers.
(Choose any person $x$.  If $x$ has met at least three of the others,
check any pair in those three have met, and so on\ldots)

\begin{figure}
\begin{framed}
\begin{theorem}
  Within every group of six people, there exists a club of three or a group of three strangers.
\end{theorem}
\begin{longFormProof}\small
  \step Let $G$ be an arbitrary group of six people.
  \comment{\ldots starts a ``for all'' block}

  \step Say that $G$ is \emph{good} if $F$ includes a club of 3 or a group of 3 strangers. 
  \comment{\ldots definition}

  \step Our goal is to prove $G$ is good.
  \comment{\ldots comment}

  \step Fix $x$ to be any one of the six people in $G$.
  \comment{\ldots definition}

  \step There two cases, each with two subcases:
  \comment{\ldots comment}

  \begin{block}[proofs: C1]
    \step[proofs: C11] \underline{Case 1}: Suppose that, among the five other than $x$,
    at least three have met $x$.
    \comment{\ldots start case block}

    \begin{block}[proofs: C1A]
      \step \underline{Subcase 1.1}: Suppose that no two among those three have met. 
      \comment{\ldots start case block}

      \step Then those three people are a group of strangers, so $G$ is good.
      \comment{\ldots statement}
    \end{block}
    \begin{block}[proofs: C1B]
      \step \underline{Subcase 1.2}: Otherwise some pair among those three have met. 
      \comment{\ldots start case block}

      \step That pair, and $x$, have all met each other, so form a club of three, so $G$ is good.
      \comment{\ldots statemen}t

    \end{block}
    \step One of the Subcases 1.1 or 1.2 must hold, and $G$ is good in each, so $G$ is good (in Case 1). 
    \comment{\ldots stmnt}
  \end{block}

  \begin{block}[proofs: C2]
    \step \underline{Case 2}:  Among the five other than $x$,
    at least three have \emph{not} met $x$.
      \comment{\ldots start case block}

    \begin{block}[proofs: C2A]
      \step \underline{Subcase 2.1}: Suppose that every pair among those three have met. 
      \comment{\ldots start case block}

      \step Those three are a club, so $G$ is good.
      \comment{\ldots statement}
    \end{block}
    
    \begin{block}[proofs: C2B]
      \step \underline{Subcase 2.2}: Otherwise some pair among those three have not met.
      \comment{\ldots start case block}
      
      \step That pair, and $x$, have all not met, so form a group of three strangers, so $G$ is good.
      \comment{\ldots statement}
    \end{block}
    \step One of the Subcases 2.1 or 2.2 must hold, and $G$ is good in each, so $G$ is good (in Case 2). 
      \comment{\ldots stmnt}
  \end{block}
  
  \step One of the Cases 1 or 2 must hold, and $G$ is good in each, so $G$ is good.
  \comment{\ldots statement}

\end{longFormProof}
\vspace*{-10pt}
\end{framed}
\vspace*{-10pt}
\caption{An example of using blocks for cases. (See the definitions in the text for context.)}\label{proofs: fig: cases}
\end{figure}

\paragraph{Remark.} 
In doing proof by cases, if the reasoning for a case
is the same as for a previously proved case,
you can just point out that the reasoning is the same,
instead of repeating it.

\paragraph{Follow each inner block by its conclusion.}
For any inner block,
the deduction that follows from the reasoning in the block
should be stated in the first line following the end of the block
(or sequence of blocks for proof by cases).
This is true for all of the four block types.
For example, in Figure~\ref{proofs: fig: for all}, just after Block 2 ends, 
Step 3 states the conclusion implied by the reasoning in the block. 
In Figure~\ref{proofs: fig: cases}, just after Blocks 6 and 7 end, 
Step 8 states the conclusion implied by the reasoning in those blocks. 

\paragraph{The main (outer) block should prove the desired fact.}
The conclusion that follows from the reasoning in the main (outer) block
should be the fact that the proof claims to prove.

\paragraph{Lemmas.}
A lemma is like a ``small'' theorem.
In these notes, we use lemmas within the context of the main proof,
to make it modular and therefore easier to understand
(just as subroutines are used to make programs more modular).
See Problem~\ref{proofs: prob: with lemma} in Section~\ref{proofs: sec: exercises} for an example.

Section~\ref{proofs: sec: exercises} has many other examples of long-form proofs for study.

% \continued

\subsection{Checklist and cheat sheet}
After you've written a proof,
put it aside for a while,
then come back with fresh eyes to check it.
Put yourself in the place of the reader,
who does not yet know anything about the line of reasoning you have in mind.
The high-level criteria is that
\emph{from reading just what you wrote,
the anticipated reader will be able to easily reconstruct and verify your entire line of reasoning.}

The main things to check are
(i) that each statement and definition is mathematically precise,
(ii) that the overall logical structure is correct,
(iii) the reasoning is complete (covers all cases),
and (iv) that there are no large gaps.
Consider each step, and make sure that if asked you could easily explain
in more detail exactly what that step says and why it holds.
Checking a proof carefully and skeptically is a learned skill.

Figure~\ref{proofs: fig: checklist} gives checklist that will help you avoid common errors.
Figure~\ref{proofs: fig: cheat sheet} gives a cheat sheet summarizing how to use the four block types.

\begin{figure}
\begin{myframed}
\begin{enumerate}
  \setlength{\labelwidth}{1em}
  \setlength{\labelsep}{1em}
  \setlength{\leftmargin}{-1em}
  \setlength{\itemsep}{0em}

\item The proof is preceded by the statement to be proved (typically as a theorem or lemma).

\item
  Each step in the proof is either
  a statement,
  definition,
  comment,
  or the start of a block.

\item
  Each definition and statement is mathematically precise and well-defined,

\item Each assumption is introduced only in the first line of some block.

\item Each block is of one of the four block types:
  contradiction, if then, for all, or cases.
\\  (A definition does not need its own block!)

\item The line following each inner block (or a series of case blocks)
  states the fact
  that follows from the reasoning in the block, according to the block's type.

\item The statement you are trying to prove follows from the main (outer) block.


\item  Each name (variable) is defined \emph{before} it is used.

\item  Each variable is referred to   only within the block in which it is defined.

\item Each statement follows directly by the justification given in that step
  from the previous steps, either by domain-specific reasoning,
  basic logic, or via one of the four block structures.
  The reader should be able to see why easily,
  without having to think very much,
  and without having to guess at what you had in mind.

\item
  The proof doesn't assume that the reader already knows anything specific about the line of reasoning.

\item
  You \emph{can} assume the reader knows basic facts about the basic structures
  you are reasoning with (e.g., numbers, trees, graphs, sets, algorithms).
  You \emph{should} use standard mathematical structures:
  numbers, tuples, sets, sequences, permutations, graphs, trees, paths, matrices, etc.

\item
  Resort to formal notation (rather than precise English) 
  only when you can't say what you need to say precisely in words alone.
  You can define a new term or notation for a concept that you introduce,
  but avoid this if possible as it increases the cognitive load on the reader.

\item Most importantly,
  from just reading what you write, the reader should be able to easily reconstruct, and hopefully verify,
  your entire line of reasoning.  

\end{enumerate}
\end{myframed}
\caption{A checklist that will help you avoid common errors in proofs.}\label{proofs: fig: checklist}
\end{figure}

\begin{figure}
\begin{multicols}{2}[
      ]\setlength{\parindent}{0in}
    \newcommand{\bx}[1]%
    {\framebox{~\begin{minipage}{0.465\textwidth}%
          ~\\#1 \vspace*{1ex} \end{minipage}~}~}

    % \bx{
    %   \underline{\emph{To show``$A$ and $B$''}}\smallskip

    %   If you have proved both ``$A$'' and ``$B$'' (separately),
    %   then you can deduce ``$A$ and $B$''.

    %   \medskip

    %   \vspace*{2pt}

    %   \underline{\emph{To use``$A$ and $B$''}}\smallskip

    %   If you have proved ``$A$ and $B$'',
    %   then you can deduce ``$A$'', and you can deduce ``$B$''.

    %   \vspace*{1pt}
    % }

    % \vspace*{2pt}

    \bx{
    %   \underline{\emph{To show``$A$ or $B$''}}\smallskip

    %   If you have proved ``$A$'', or you have proved ``$B$'',
    %   then you can deduce ``$A$ or $B$''.

    %   \medskip

    %   \vspace*{2pt}

      % \underline{\emph{To use``$A$ or $B$'' (proof by cases)}}\smallskip 

    \underline{\emph{Proof by cases, given that ``$A$ or $B$'' holds}}\smallskip 

      \begin{steps}

      \item[1.] \ldots
        
        \begin{steps}
          
        \item[2.1.] \underline{Case 1.}  Assume $A$. 

        \item[2.2.] \emph{ \ldots~within this scope, $A$ holds \ldots}

        \item[] ~~~~~~\emph{\ldots~somehow prove $C$ \ldots }

        \item[2.3.] $C$
          
        \end{steps}

        \begin{steps}
          
        \item[3.1.] \underline{Case 2.}  Assume $B$. 

        \item[3.2.] \emph{ \ldots~within this scope, $B$ holds \ldots}

        \item[] ~~~~~~\emph{\ldots~somehow prove $C$ \ldots }

        \item[3.3.] $C$
          
        \end{steps}

      \item[4.] By Cases 1 and 2, $C$ holds.

      \end{steps}
    }

    \vspace*{2pt}

    \bx{
      % \underline{\emph{To use``if $A$ then $B$''}}\smallskip

      % If you have proved ``if $A$ then $B$'', and you have proved ``$A$'',
      % then you can deduce ``$B$''.

      % \medskip

      % \vspace*{2pt}

      \underline{\emph{``If  then'' blocks}}\smallskip
      \begin{steps}

      \item[1.] \ldots
        
        \begin{steps}
          
        \item[2.1.] Assume $A$. 

        \item[2.2.] \emph{ \ldots~within this scope, $A$ holds \ldots}

        \item[] ~~~~~~\emph{\ldots~somehow prove $B$ \ldots }
          
        \item[2.3.] B
          
        \end{steps}
        
      \item[3.] By Block 2, if $A$, then $B$.
      \end{steps}

      \vspace*{1pt}

      \dotfill

      \emph{Alternatively, show the contrapositive:\\\hspace*{3em} ``if not $B$, then not $A$''.}

      \vspace*{1pt}
    }

    \newpage
    
    % \bx{
    %   \underline{\emph{To show ``\,$\exists x\in S.~P(x)$''}}

    %   Prove that $P(y)$ holds for a particular $y\in S$.

    %   \medskip

    %   \underline{\emph{To use ``\,$\exists x\in S.~P(x)$''}}

    %   \begin{steps}

    %   \item Let $y$ be an element of $S$ such that $P(y)$.

    %   \item \emph{ \ldots~Now you can use $y$~\ldots}
    %   \item[] \emph{\ldots~~~and that $y\in S$ and $P(y)$~~\ldots}

    %   \end{steps}
    % }

    \bx{
      % \underline{\emph{To use ``not $A$''}}

      % Prove ``$A$ and not $A$'' to reach a \textbf{contradiction}.


      % \medskip

      \underline{\emph{Proof by contradiction}}

      \begin{steps}

      \item[1.] \ldots

        \begin{steps}
          
        \item[2.1.] Assume \textbf{for contradiction} that $B$ holds. 


        \item[2.2.] \emph{ \ldots~within this scope, $B$ holds \ldots}

        \item[] ~~~~~~\emph{\ldots~prove any contradiction \ldots }

        \item[2.3.] contradiction
          
        \end{steps}
        
      \item[3.] By Block 2, $B$ is false.
      \end{steps}
    }

    \vspace*{2pt}


    \bx{
      % \underline{\emph{To use ``\,$\forall x\in S.~P(x)$''}}

      % If you have proved ``\,$\forall x\in S.~P(x)$'',
      % and you know $y\in S$ for some particular $y$,
      % you can deduce $P(y)$ for that $y$.

      % \medskip 

      \underline{\emph{``For all''  blocks}}

      \begin{steps}

      \item[1.]\ldots

        \begin{steps}
          
        \item[2.1.] Let $x$ be an \textbf{arbitrary} element of $S$.~\footnote
          {The definition in Step 2.1 assumes $S$ is non-empty.
            This assumption is implicitly introduced by the block.
            For other types of definitions, blocks are not needed.}

        \item[2.2.] \emph{ \ldots~within this scope, $x$ is defined \ldots}

        \item[] ~~~~~~\emph{\ldots~somehow prove $P(x)$ \ldots }

        \item[2.3.] $P(x)$
          
        \end{steps}
        
      \item[3.] By Block 2, every $x$ in $S$ has property $P(x)$.
      \end{steps}
    }

    \vspace*{2pt}
    
    \bx{
      \underline{\emph{Basic logic:}}  (``$\equiv$'' means ``is equivalent to'')

      \begin{tabular}{r@{~}l}
        ``if $A$ then $B$'' &  $\equiv$ ``if not $B$ then not $A$'' 
        \\
        ``if $A$ then $B$''  & $\equiv$ ``$B$ or not $A$''
        \\
        ``not ($A$ and $B$)'' &  $\equiv$ ``(not $A$) or (not $B$)'' 
        \\
        ``not ($A$ or $B$)''  & $\equiv$ ``(not $A$) and (not $B$)'' 
        \\
        ``not (not $A$)'' & $\equiv$ ``$A$''
        \\
        ``not $\forall x\in S.~P(x)$'' & $\equiv$ ``$\exists x\in S.~\text{not } P(x)$'' 
        \\
        ``not $\exists x\in S.~P(x)$'' & $\equiv$ ``$\forall x\in S.~\text{not } P(x)$'' 
      \end{tabular}
    }
  \end{multicols}

  \paragraph{Each block introduces an assumption that lasts for the scope of the block.}
  The general principle is that
  \emph{whenever we want to explore the consequences of some assumption,
    we begin a new block, with a step that makes that assumption}.
  The block delimits the scope of that assumption.
  There are four block types:
  ``proof by contradiction'',  ``if then'',  ``for all'',  and ``cases''.
  (For the ``for all'' block, what is the underlying assumption that necessitates a block?
  That $S$ is not empty!)
\caption{A cheat sheet summarizing the four block types and how to use them.}\label{proofs: fig: cheat sheet}
\end{figure}

%%% Local Variables:
%%% mode: latex
%%% TeX-master: "long_form_proofs"
%%% End:
